\section{The financial retention problem}
\label{sec:introduction}

A portfolio research team accumulates strategies that it must test, monitor,
document, and maintain. Some are used in portfolios today. Others remain in the
library because their signals, filters, allocation rules, or risk controls may
be useful in later research. Maintaining every strategy consumes resources;
removing a strategy can reduce that burden. The manager's problem is to decide
which strategies can leave the active library without giving up either current
operating choices or capabilities needed to build future strategies.

Consider two strategies that offer the same current portfolio opportunity.
One also carries a capability needed by a possible new strategy. If the team
retains only the other strategy, today's opportunity is preserved but the new
strategy may become unavailable. A performance-based redundancy check cannot
detect that loss. Conversely, keeping every strategy that carries a useful
capability can be unnecessarily expensive when several entries preserve the
same capabilities. The decision therefore concerns the retained library as a
whole.

Here, removal means leaving the active maintenance and validation set. Legal
records, source code, research history, and mandatory archives remain subject
to their own retention rules. A capability is an explicitly declared input to
the modeled research process. The model applies when retaining its designated
carriers is what makes that capability available; it does not assume that a
strategy label contains all of an institution's knowledge.

We represent the manager's two preservation targets in simple terms. The
\emph{operating frontier} records the best current operating value offered by
the source library in each modeled condition. The \emph{generative closure}
records the capabilities available after applying the declared rules for
combining the retained strategies' components. Innovation-safe compression
finds a least-burden subset of the source library that preserves both targets.
The burden can represent a specified validation or governance index; its
interpretation must be fixed before looking at subsequent investment outcomes.

The contribution is a conditional preservation principle for a maintained
strategy library, an exact optimization representation of that principle, and
counterexamples showing why weaker deletion rules can fail. The financial
question concerns the availability of reusable inputs to subsequent research.
The covering algorithms used to implement the retention decision are
established methods. Our financial evidence measures the declared burden and
research opportunities separately from subsequent investment performance.

The mathematical argument follows that decision. First, we specify when the
two targets summarize the library's role in a research process. We prove that
preserving them then preserves the process's operating rewards and future
choices. An identity-dependent project supplies a counterexample when the
assumption fails. Next, for a model in which capabilities are simply the union
of retained components, we show that the decision is a weighted covering
problem. This connects the financial specification to established optimization
methods \citep{Karp1972,Chvatal1979}. A separate counterexample proves that
repeatedly deleting the heaviest redundant entry can leave an arbitrarily
expensive library even when every deletion is safe.

The computational study checks exact agreement and the burdens returned by
several methods on a registered synthetic design. The financial application
then examines portfolio-strategy libraries formed from point-in-time common
equities and ETFs. It finds substantial reductions in a validation-burden
index while preserving more research options than a budget-matched library
that preserves only current operation. The protected policy exercises none of
the additional common-equity options; its single ETF replacement loses
held-out value. A calibrated synthetic experiment separately shows why an
available positive-value option can remain too difficult to recognize at a
small signal-to-noise ratio.

These results distinguish preserving an opportunity, deciding to use it, and
realizing value from it. They establish a conditional retention principle and
bounded computational and financial findings. The financial application is
retrospective development evidence, and its burden index is not a measured
institutional cost saving. The assumptions and the counterexamples below
state precisely where the principle applies.
